\scanpage{99}
\mspage{98}
\margin{\diagram{\begin{tikzcd}[ampersand replacement=\&, column sep=small,row sep=small]
\& P \arrow[d,no head] \& U \arrow[d] \\
E_0 \arrow[r] \& P/\mu=P_0 \arrow[d] \& U_0=U/\mu \\
\& \unsure{T} \&
\end{tikzcd}}\quad\unsure{tor.}}
$1^\circ)$ Une structure de $\mathcal E$-groupoïde (central)
\inserted{$C_0$}
sur $E_0$ comme ens d'objets
\deleted{\unsure{$2^\circ)$}} $\unsure{C}$ fournit un paradigme de
$P_{\unsure{0|q_0}}\longrightarrow T$ comme $U_0$-torseur,
et $E_0\mathbin{\text{---}}P_0$ pour marquer des
pts pour groupoïde fondamental.

\margin{\diagram{\begin{tikzcd}[ampersand replacement=\&, column sep=small,row sep=small]
1 \arrow[r,no head] \& \mathcal E \arrow[r] \arrow[d] \&
\underset{\unsure{\pi_1^{\prime\prime}(P_0)}}{G_0} \arrow[r] \arrow[dl,"\varphi"] \&
\underset{\unsure{\pi_1^{\prime\prime}(T)}}{G} \arrow[r,no head] \&1 \\
\& \mu \&\&\&
\end{tikzcd}}}
$2^\circ)$ Un \deleted{$\mu\longrightarrow$} $\mu$-rev. principal
\deleted{\unsure{$\widetilde E|E$}} $E$
de $\unsure{E_0|P_0}$, défini par un système
local sur $C_0$ sur $P_0$, avec op. de
translation de $\mu$ dessus (de façon
que $C_0\simeq C/\mu$), l'homomorphisme
\[
\pi_1(C_0,x_0)\longrightarrow\mu\simeq\operatorname{Aut}(\unsure{E|\mathcal E}_{x_0})
\]
prolongeant
\inserted{$\mathcal E\hookrightarrow\pi_1(C_0,x_0)$}
l'homomorphisme de passage
au quotient $\mathcal E\mathbin{\text{---}}\mu$
\margin{i.e. on a\quad $\pi_{1\mathrm{ab}}(C_0)\longrightarrow\mu$\quad $\varphi_0$}
\margin{\deleted{$\operatorname{Ker}\varphi$}
$1\longrightarrow\pi\longrightarrow\operatorname{Ker}\varphi
\xrightarrow{\uuline{\mathrm{surj}}}G\mathbin{\text{---}}1$}

\diagram{\begin{tikzcd}[ampersand replacement=\&]P\arrow[d,"f"']\\T\end{tikzcd}}
\deleted{fibré} torseur sur le fibré
en groupes
top. connexes $U_T$
(syst. loc. des $\pi_1(U_t)$) est $\pi_T$
\diagram{\begin{tikzcd}[ampersand replacement=\&]\Pi_1P\arrow[d]\\\Pi_1T\end{tikzcd}}
$f^*(\pi_T)$ \deleted{\unsure{hom}} opère sur \deleted{\unsure{$\pi$}} $\Pi_1P$
\[
f^*(\pi_T)\longrightarrow\operatorname{Int}\Pi_1P
\]
$\gamma_T$ un sous-fibré en groupes inv. de $U_T$
(\unsure{noté comm.})
\[
U_{0T}=U_T/\gamma_T
\qquad \text{Syst. loc. des }\pi_1\text{ est }\mathcal E_T
\]
\[
0\longrightarrow\pi_T\mathbin{\text{---}}\mathcal E_T
\longrightarrow\gamma_T\longrightarrow1
\]
